% !TEX root = ../main.tex

\section{Requirement Review}
\label{sec:results-requirements}

\scaffold{
  One sentence introducing \cref{tab:requirements}, which checks the proposal \cite{benson2026proposal} and the evaluation plan of milestone~1 against the results; no further prose.
}

\begin{table}[htbp]
  \centering
  \footnotesize
  \titledcaption{Requirements and Results}{\scaffoldinline{P: proposal, M: evaluation plan of milestone~1; status met, partly met or not met.}}
  \label{tab:requirements}
  \begin{tabularx}{\linewidth}{@{}l >{\raggedright\arraybackslash}p{0.34\linewidth} >{\raggedright\arraybackslash}X l@{}}
    \toprule
    & \textbf{Requirement} & \textbf{Result} & \textbf{Status} \\
    \midrule
    P & Detect presence and wiring automatically & all available pedal types and both wirings & met \\
    P & Real time, no perceptible latency & 188 updates/s, recognition \qty{0.1}{\second} (\qty{0.3}{\second} worst case) & met \\
    P & Value linear to the pedal angle & ratiometric wiper reading, drive curve; linearity not measured & partly \\
    P & One to four TRS jacks, class-compliant USB MIDI & four jacks, USB MIDI 1.0 & met \\
    P & Calibration and filtering of the channels & rail measurement, MIDI hysteresis & met \\
    P & Optional DIN MIDI output & -- & not met \\
    M & Supply, SPI, self-test & \cref{sec:results-board} & met \\
    M & Known networks, accuracy, linearity & real pedals and ohmmeter only & partly \\
    M & End-to-end latency & estimated, not measured & partly \\
    \bottomrule
  \end{tabularx}
\end{table}
